\lecture{22}{20. November 2025}{Combined Loading}

\begin{problemwithimage}{f13_5}{13-5}
  The cylindrical tank is fabricated by welding a strip of thin plate helically, making an angle $\theta$ with the longitudinal axis of the tank.
  If the strip has a width $w$ and thickness $t$, and the gas within the tank of diameter $d$ is subjected to a pressure $p$, show that the normal stress developed along the strip is given by $\sigma_{\theta} = \left( p d / 8 t\right) \left( 3 - \cos 2\theta \right)$
\end{problemwithimage}

\begin{derivation}
  We have the general results, that the longitudinal stress is:
  \begin{equation*}
    \sigma_x = \frac{p d}{4t}
  \end{equation*}
  and the hoop stress is
  \begin{equation*}
    \sigma_{\phi} = \frac{p d}{2t}
  \end{equation*}
  whilst the torsion in the vessel is zero in the ideal case:
  \begin{equation*}
    \tau_{x\phi} = 0
  .\end{equation*}
  
  The strip is rotated at an angle $\theta$ to the hoop.
  Here we can use the stress-transformation formula:
  \begin{align*}
    \sigma_{\theta} &= \sigma_x \cos^2 \theta + \sigma_{\phi} \sin^2 \theta + 2 \tau_{x \phi} \cos \theta \sin \theta \\
    &= \sigma_x \cos^2 \theta + \sigma_{\phi} \sin^2\theta \\
    &= \frac{p d}{4t} \cos^2 \theta + \frac{p d}{2t} \sin^2 \theta \\
    &= \frac{p d}{4t} \left( \cos^2 \theta + 2 \sin^2 \theta \right) \\
    &= \frac{p d}{4t} \left( 1 + \sin^2 \theta \right) \\
    1 + \sin^2 \theta &= \frac{3 - \cos 2\theta}{2} \\
    \sigma_{\theta} &= \frac{p d}{4t} \cdot \frac{3 - \cos 2\theta}{2} \\
    &= \frac{p d}{8t} \cdot \left( 3 - \cos 2\theta \right)
  .\end{align*}
\end{derivation}

\finalresult{
  \sigma_{\theta} = \frac{p d}{8t} \cdot \left( 3 - \cos 2\theta \right)
}


\begin{problemwithimage}{f13_11}{13-11}
  The staves or vertical members of the wooden tank are held together using semicircular hoops having a thickness of \qty{12}{\mm} and a width of \qty{50}{\mm}.
  Determine the normal stress in hoop $AB$ if the tank is subjected to an internal gauge pressure of \qty{14}{\kPa} and this loading is transmitted directly to the hoops.
  Also, if \qty{6}{\mm} diameter bolts are used to connect each hoop together, determine the tensile stress in each bolt at $A$ and $B$.
  Assume hoop $AB$ supports the pressure loading within a \qty{300}{\mm} length of the tank as shown.
\end{problemwithimage}

\begin{derivation}
  We have a ring-shaped slice of length $L = \qty{300}{\mm}$, diameter $d = \qty{450}{\mm}$ and an internal gauge pressure of $p = \qty{14}{\kPa} $.

  We now split the barrel in half and consider equilibrium for one half.
  Here the projected area is:
  \begin{equation*}
    A_{\mathrm{proj}} = d L = \qty{300}{\mm} \cdot \qty{450}{\mm} = \qty{135e3}{\mm^2} 
  .\end{equation*}
  The pressure force on this side is just the pressure times the area:
  \begin{equation*}
    F_p = p \cdot A_{\mathrm{proj}} = \qty{14}{\kPa} \cdot \qty{135e3}{\mm^2} = \qty{1890}{\N} 
  .\end{equation*}
  The only thing resisting this pressure force is the hoop tension.
  As the hoop pulls both on the front and backside of the barrel the tension in the hoop must be:
  \begin{align*}
    \sum F_x &= 0 \\
    F_p - 2 T &= 0 \\
    T &= \frac{F_p}{2} \\
      &= \frac{\qty{1890}{\N} }{2} = \qty{945}{\N}
  .\end{align*}
  
  The stress in the hoop is then:
  \begin{equation*}
    \sigma_{\mathrm{hoop}} = \frac{T}{A} = \frac{\qty{945}{\N}}{\qty{12}{\mm} \cdot \qty{50}{\mm} } = \qty{1,575}{\MPa}
  .\end{equation*}
  
  The bolts at $A$ and $B$ must transmit the entirety of the hoop force.
  This means:
  \begin{equation*}
    F_{\mathrm{bolt}} = T
  .\end{equation*}
  Which means the stress in the bolts are:
  \begin{equation*}
    \sigma_{\mathrm{bolt}} = \frac{F_{\mathrm{bolt}}}{\pi \cdot \frac{d_{\mathrm{bolt}}^2}{4}} = \frac{\qty{945}{\N} }{\pi \cdot \frac{\left( \qty{6}{\mm}  \right)^2}{4}} = \qty{33,4}{\MPa}
  .\end{equation*}
  Both of these are in tension.
\end{derivation}

\finalresult{
\begin{aligned}
  \sigma_{\mathrm{hoop}} &= \qty{1,575}{\MPa} \\
  \sigma_{\mathrm{bolt}} &= \qty{33,4}{\MPa} 
.\end{aligned}
}
